\begin{Example}
\label{ex:sturm-liouville-robin-example}
Consider the Sturm-Liouville eigenvalue problem consisting of the differential equation
\begin{equation}\label{eq:sturm-liouville-robin-example-ode}
-u''-\lambda u=0
\end{equation}
and the boundary conditions
\begin{equation}\label{eq:sturm-liouville-robin-example-left-condition}
u(0)=0,
\end{equation}

\begin{equation}\label{eq:sturm-liouville-robin-example-right-condition}
\cos\beta\,u(1)+\sin\beta\,u'(1)=0.
\end{equation}
Without loss of generality, we assume that $-\pi/2<\beta\leq\pi/2$.

This equation arises from (for instance) the separation of variables of a one-dimensional heat conduction problem modeling a long, thin rod that is insulated on the sides and with one end held at a fixed temperature and the other radiating heat.

Depending on the sign of $\lambda$ we get the following two-parameter families of real solutions for \eqref{eq:sturm-liouville-robin-example-ode}:
\begin{equation}\label{eq:sturm-liouville-robin-example-general-solutions}
y(x)=\begin{cases}
A\cosh\sqrt{-\lambda}x+B\sinh\sqrt{-\lambda}x,&\lambda<0,\\
A+Bx,&\lambda=0,\\
A\cos\sqrt{\lambda}x+B\sin\sqrt{\lambda}x,&\lambda>0.
\end{cases}
\end{equation}
Applying the boundary condition \eqref{eq:sturm-liouville-robin-example-left-condition} implies
\begin{equation}\label{eq:sturm-liouville-robin-example-left-reduced-solutions}
y(x)=\begin{cases}
B\sinh\sqrt{-\lambda}x,&\lambda<0,\\
Bx,&\lambda=0,\\
B\sin\sqrt{\lambda}x,&\lambda>0.
\end{cases}
\end{equation}
We now analyze the second boundary condition.\begin{enumerate}\item For $\lambda<0$ we let $\omega=\sqrt{-\lambda}$. Then \eqref{eq:sturm-liouville-robin-example-right-condition} becomes
\begin{equation}\label{eq:sturm-liouville-robin-example-negative-boundary-equation}
B(\cos\beta\sinh\omega+\omega\sin\beta\cosh\omega)=0.
\end{equation}
\begin{figure}[H]
\centering
\includegraphics[width=.78\linewidth]{figures/sturm-liouville-negative-spectrum-source.png}
\caption{Solutions of $\tanh\omega=-\omega\tan\beta$.}\label{fig:sturm-liouville-negative-roots}
\end{figure}Note that if $\beta=\pi/2$, there is no solution. Thus, there is a nontrivial solution if and only if there is an $\omega>0$ satisfying
\begin{equation}\label{eq:sturm-liouville-robin-example-negative-root-equation}
\tanh\omega=-\omega\tan\beta.
\end{equation}
We see from Figure~\ref{fig:sturm-liouville-negative-roots} that if $\tan\beta\in(-1,0)$ (i.e., $\beta\in(-\pi/4,0)$), there is exactly one solution $\omega_0$, otherwise there is no solution. If a solution exists, it corresponds to the eigenvalue $\lambda_0=-\omega_0^2$ and the eigenfunction $u_0(x)=\sinh\omega_0x$.\item For $\lambda=0$, \eqref{eq:sturm-liouville-robin-example-right-condition} becomes
\begin{equation}\label{eq:sturm-liouville-robin-example-zero-boundary-equation}
B(\cos\beta+\sin\beta)=0.
\end{equation}
Thus, there is a nontrivial solution if and only if $\beta=-\pi/4$. In this case $\lambda_0=0$ is an eigenvalue and $u_0(x)=x$ is the corresponding eigenfunction.\item For $\lambda>0$ we let $\omega=\sqrt{\lambda}$. Then \eqref{eq:sturm-liouville-robin-example-right-condition} becomes
\begin{equation}\label{eq:sturm-liouville-robin-example-positive-boundary-equation}
B(\cos\beta\sin\omega+\omega\sin\beta\cos\omega)=0.
\end{equation}
For $\beta\in(-\pi/2,\pi/2)$ we have nontrivial solutions if and only if there is an $\omega>0$ which satisfies
\begin{equation}\label{eq:sturm-liouville-robin-example-positive-root-equation}
\tan\omega=-\omega\tan\beta.
\end{equation}
\begin{figure}[H]
\centering
\includegraphics[width=.78\linewidth]{figures/sturm-liouville-positive-spectrum-source.png}
\caption{Solutions of $\tan\omega=-\omega\tan\beta$. (a) $\tan\beta>0$. (b) $\tan\beta<0$.}\label{fig:sturm-liouville-positive-roots}
\end{figure}From Figure~\ref{fig:sturm-liouville-positive-roots} we see that for any $\beta\in(-\pi/2,\pi/2)$ this equation has an infinite family of solutions $\omega_n$, $n=1,2,3,\ldots$ such that $\lim_{n\to\infty}\omega_n=\infty$. Thus there is an infinite family of eigenvalues $\lambda_n=\omega_n^2$ with corresponding eigenfunctions $u_n(x)=\sin\omega_nx$, $n=1,2,3,\ldots$.

If $\beta=\pi/2$, we solve \eqref{eq:sturm-liouville-robin-example-positive-boundary-equation} directly to get the eigenvalues $\lambda_n=(2n-1)^2\pi^2/4$ with corresponding eigenfunctions $\sin[(2n-1)\pi x/2]$, $n=1,2,3,\ldots$.\end{enumerate}
\end{Example}
